\section{Background \label{theory}}

This chapter provides the theoretical foundations for the methods used in this thesis.
Section~\ref{sec:chaotic} introduces chaotic dynamical systems and the three benchmark datasets.
Section~\ref{sec:rnn} reviews recurrent neural networks and the vanishing gradient problem.
Section~\ref{sec:lstm} presents the Long Short-Term Memory architecture.
Section~\ref{sec:esn} covers Reservoir Computing and Echo-State Networks.
Section~\ref{sec:related} discusses related work and positions this thesis within the literature.

% ============================================================================
\subsection{Chaotic Dynamical Systems}
\label{sec:chaotic}
% ============================================================================

A dynamical system is called \emph{chaotic} if it is deterministic yet exhibits sensitive dependence on initial conditions: two trajectories starting from arbitrarily close initial states diverge exponentially over time \cite{strogatz2018nonlinear}. This sensitivity, popularly known as the \emph{butterfly effect}, makes long-term prediction inherently challenging even when the governing equations are fully known.

Chaos is quantified through \emph{Lyapunov exponents}, which measure the average rate of exponential divergence of nearby trajectories. For a system with state vector $\mathbf{x} \in \mathbb{R}^n$, the maximal Lyapunov exponent $\lambda_{\max}$ is defined as
\begin{equation}
    \lambda_{\max} = \lim_{t \to \infty} \frac{1}{t} \ln \frac{\| \delta\mathbf{x}(t) \|}{\| \delta\mathbf{x}(0) \|},
    \label{eq:lyapunov}
\end{equation}
where $\delta\mathbf{x}(t)$ denotes an infinitesimal perturbation at time $t$ \cite{wolf1985determining}. A positive maximal Lyapunov exponent ($\lambda_{\max} > 0$) indicates chaotic behavior. Chaotic systems are bounded---their trajectories remain in a finite region of phase space called a \emph{strange attractor}---yet never exactly repeat.

The prediction horizon of a chaotic system scales inversely with $\lambda_{\max}$: the \emph{Lyapunov time} $T_\lambda = 1/\lambda_{\max}$ gives a characteristic timescale beyond which predictions lose reliability. For data-driven models, the goal is to achieve accurate predictions up to or beyond this timescale.

The following three benchmark systems are used throughout this thesis. They span a range of complexity: a continuous-time ordinary differential equation (Lorenz), a delay differential equation (Mackey-Glass), and a real-world experimental time series (Santa Fe laser).

% --------------------------------------------------------------------------
\subsubsection{The Lorenz System}
\label{sec:lorenz}
% --------------------------------------------------------------------------

The Lorenz system, introduced by Edward Lorenz in 1963 \cite{lorenz1963deterministic}, is one of the earliest and most studied examples of deterministic chaos. Originally derived as a simplified model of atmospheric convection, it is governed by the following system of ordinary differential equations:
\begin{align}
    \dot{x} &= \sigma(y - x), \label{eq:lorenz_x} \\
    \dot{y} &= x(\rho - z) - y, \label{eq:lorenz_y} \\
    \dot{z} &= xy - \beta z, \label{eq:lorenz_z}
\end{align}
where $\sigma$, $\rho$, and $\beta$ are positive parameters. The canonical parameter values $\sigma = 10$, $\rho = 28$, and $\beta = 8/3$ yield chaotic behavior. The system possesses three equilibrium points: the origin and two symmetric points $C^\pm = (\pm\sqrt{\beta(\rho-1)}, \pm\sqrt{\beta(\rho-1)}, \rho-1)$, both of which are unstable.

The Lorenz attractor has a characteristic butterfly-shaped double-scroll structure, with the trajectory alternating unpredictably between the two lobes. The maximal Lyapunov exponent at the canonical parameters is approximately $\lambda_{\max} \approx 0.9$, giving a Lyapunov time of about $T_\lambda \approx 1.1$ time units.

In this thesis, Lorenz trajectories are generated by numerical integration using \texttt{scipy.integrate.solve\_ivp} with a time step of $\Delta t = 0.01$ and initial condition $\mathbf{x}_0 = (1, 1, 1)$. An initial transient is discarded to ensure the trajectory has settled onto the attractor. The $x$-component is used as the primary univariate prediction target.

% --------------------------------------------------------------------------
\subsubsection{The Mackey-Glass System}
\label{sec:mackeyglass}
% --------------------------------------------------------------------------

The Mackey-Glass equation, proposed in 1977 \cite{mackeyglass1977}, models physiological feedback systems such as the regulation of white blood cell production. It is a delay differential equation (DDE) of the form
\begin{equation}
    \frac{dx}{dt} = \frac{\beta_{\text{MG}} \, x_\tau}{1 + x_\tau^n} - \gamma \, x,
    \label{eq:mackeyglass}
\end{equation}
where $x_\tau = x(t - \tau)$ denotes the delayed state, $\tau$ is the time delay, $\beta_{\text{MG}}$ and $\gamma$ are rate constants, and $n$ controls the nonlinearity. The standard parameter values used in this thesis are $\beta_{\text{MG}} = 0.2$, $\gamma = 0.1$, $n = 10$, and $\tau = 17$.

The qualitative behavior of the system depends critically on the delay $\tau$: for $\tau < 4.53$, the system converges to a stable fixed point; for intermediate delays, it exhibits periodic oscillations; and for $\tau = 17$, the system enters a chaotic regime. The resulting time series has a rich spectral structure that makes it a challenging benchmark for prediction algorithms.

The Mackey-Glass system has been extensively used as a benchmark since the Santa Fe Time Series Competition \cite{weigend1994time}. Its moderate complexity---more challenging than low-dimensional ODEs but with well-understood dynamics---makes it an ideal intermediate test case.

In this thesis, the Mackey-Glass equation is integrated numerically with parameters as specified above, and the resulting scalar time series $x(t)$ is sampled at discrete intervals.

% --------------------------------------------------------------------------
\subsubsection{The Santa Fe Laser Dataset}
\label{sec:santafe}
% --------------------------------------------------------------------------

The Santa Fe laser dataset originates from the Santa Fe Time Series Prediction Competition held in 1991--1993 \cite{weigend1994time}. It consists of experimental intensity measurements from an NH\textsubscript{3} far-infrared laser operating in a chaotic regime \cite{hubner1989dimensions}. Unlike the Lorenz and Mackey-Glass systems, this dataset is purely empirical: the underlying dynamics are high-dimensional, the measurement contains instrumental noise, and no simple closed-form equations are available.

The dataset (labeled ``Dataset A'' in the competition) contains 1\,000 data points of laser intensity measurements. The time series exhibits collapse events---episodes where the intensity drops sharply before recovering---interspersed with intervals of approximately periodic oscillation. This mixture of regularity and unpredictability, combined with measurement noise, makes the Santa Fe laser dataset a challenging real-world benchmark.

Including this experimental dataset alongside the synthetic Lorenz and Mackey-Glass systems allows us to assess whether conclusions drawn from controlled simulations transfer to real-world settings with unknown noise characteristics and higher-dimensional latent dynamics.

% ============================================================================
\subsection{Recurrent Neural Networks}
\label{sec:rnn}
% ============================================================================

Recurrent neural networks (RNNs) are a class of neural networks designed for processing sequential data. Unlike feedforward networks, RNNs maintain a hidden state $\mathbf{h}_t \in \mathbb{R}^h$ that is updated at each time step based on the current input $\mathbf{x}_t \in \mathbb{R}^d$ and the previous hidden state:
\begin{equation}
    \mathbf{h}_t = f(\mathbf{W}_{hh} \mathbf{h}_{t-1} + \mathbf{W}_{xh} \mathbf{x}_t + \mathbf{b}_h),
    \label{eq:rnn}
\end{equation}
where $\mathbf{W}_{hh} \in \mathbb{R}^{h \times h}$ and $\mathbf{W}_{xh} \in \mathbb{R}^{h \times d}$ are weight matrices, $\mathbf{b}_h$ is a bias vector, and $f(\cdot)$ is a nonlinear activation function (typically $\tanh$). The output $\mathbf{y}_t$ is computed from the hidden state through an output layer:
\begin{equation}
    \mathbf{y}_t = \mathbf{W}_{hy} \mathbf{h}_t + \mathbf{b}_y.
    \label{eq:rnn_output}
\end{equation}

RNNs are trained using \emph{Backpropagation Through Time} (BPTT), which unfolds the recurrence over the input sequence and applies standard backpropagation. However, BPTT suffers from the \emph{vanishing and exploding gradient problem} \cite{bengio1994learning, pascanu2013difficulty}: when gradients are propagated backward through many time steps, they either shrink exponentially (vanishing) or grow uncontrollably (exploding).

To see why, consider the gradient of the loss $\mathcal{L}$ with respect to the hidden state at time step $k < t$:
\begin{equation}
    \frac{\partial \mathcal{L}}{\partial \mathbf{h}_k} = \frac{\partial \mathcal{L}}{\partial \mathbf{h}_t} \prod_{i=k+1}^{t} \frac{\partial \mathbf{h}_i}{\partial \mathbf{h}_{i-1}}.
    \label{eq:vanishing}
\end{equation}
Each factor $\partial \mathbf{h}_i / \partial \mathbf{h}_{i-1}$ involves the recurrent weight matrix $\mathbf{W}_{hh}$ and the derivative of the activation function. If the spectral radius of $\mathbf{W}_{hh}$ is less than one, the product of Jacobians shrinks exponentially with $t - k$, causing vanishing gradients. If it exceeds one, gradients explode. In practice, this means that standard RNNs cannot effectively learn dependencies spanning more than approximately 10--20 time steps \cite{bengio1994learning}.

% ============================================================================
\subsection{Long Short-Term Memory Networks}
\label{sec:lstm}
% ============================================================================

Long Short-Term Memory (LSTM) networks were introduced by Hochreiter and Schmidhuber \cite{hochreiter1997long} specifically to address the vanishing gradient problem. The key innovation is a \emph{cell state} $\mathbf{c}_t$ that acts as a memory highway, allowing information to flow across many time steps with minimal degradation. The flow of information into and out of the cell state is controlled by three multiplicative gates.

\subsubsection{Architecture and Gate Equations}

At each time step $t$, given input $\mathbf{x}_t \in \mathbb{R}^d$ and previous hidden state $\mathbf{h}_{t-1} \in \mathbb{R}^h$, the LSTM computes:

\textbf{Forget gate} -- decides what information to discard from the cell state:
\begin{equation}
    \mathbf{f}_t = \sigma(\mathbf{W}_f [\mathbf{h}_{t-1}, \mathbf{x}_t] + \mathbf{b}_f),
    \label{eq:forget_gate}
\end{equation}

\textbf{Input gate} -- decides what new information to store:
\begin{equation}
    \mathbf{i}_t = \sigma(\mathbf{W}_i [\mathbf{h}_{t-1}, \mathbf{x}_t] + \mathbf{b}_i),
    \label{eq:input_gate}
\end{equation}

\textbf{Candidate cell state}:
\begin{equation}
    \tilde{\mathbf{c}}_t = \tanh(\mathbf{W}_c [\mathbf{h}_{t-1}, \mathbf{x}_t] + \mathbf{b}_c),
    \label{eq:candidate}
\end{equation}

\textbf{Cell state update}:
\begin{equation}
    \mathbf{c}_t = \mathbf{f}_t \odot \mathbf{c}_{t-1} + \mathbf{i}_t \odot \tilde{\mathbf{c}}_t,
    \label{eq:cell_update}
\end{equation}

\textbf{Output gate} -- decides what to output:
\begin{equation}
    \mathbf{o}_t = \sigma(\mathbf{W}_o [\mathbf{h}_{t-1}, \mathbf{x}_t] + \mathbf{b}_o),
    \label{eq:output_gate}
\end{equation}

\textbf{Hidden state}:
\begin{equation}
    \mathbf{h}_t = \mathbf{o}_t \odot \tanh(\mathbf{c}_t),
    \label{eq:hidden}
\end{equation}
where $\sigma(\cdot)$ denotes the sigmoid function, $\odot$ denotes element-wise multiplication, $[\cdot, \cdot]$ denotes concatenation, and the $\mathbf{W}$ matrices and $\mathbf{b}$ vectors are learnable parameters.

The forget gate $\mathbf{f}_t$, introduced by Gers et al. \cite{gers2000learning}, enables the network to selectively erase parts of the cell state. When $\mathbf{f}_t \approx 1$, information is preserved; when $\mathbf{f}_t \approx 0$, it is discarded. This gating mechanism creates paths through the network along which gradients can flow without vanishing, enabling learning of long-range dependencies.

This thesis uses the \emph{vanilla LSTM} architecture---with input gate, forget gate, output gate, and cell state, but without peephole connections or coupled input-forget gates. Greff et al. \cite{greff2017lstm} conducted a comprehensive comparison of eight LSTM variants and concluded that none of the more complex modifications (peephole connections, coupled gates, full gate recurrence) significantly improve performance over the vanilla architecture. This finding justifies the simpler design used here, which also simplifies the parameter count analysis.

\subsubsection{Parameter Count}

For a single-layer LSTM with hidden size $h$ and input dimension $d$, each of the four gate computations (forget, input, candidate, output) requires a weight matrix of size $h \times (h + d)$ and a bias vector of size $h$. The total number of parameters for one LSTM layer is therefore:
\begin{equation}
    N_{\text{LSTM}} = 4 \left( h^2 + h \cdot d + h \right) = 4h(h + d + 1).
    \label{eq:lstm_params}
\end{equation}
Including the final linear output layer from $h$ to the output dimension $d_{\text{out}}$ adds $h \cdot d_{\text{out}} + d_{\text{out}}$ parameters. All parameters in an LSTM are trainable, meaning they are optimized through gradient-based backpropagation.

\subsubsection{Architecture Design for Time-Series Prediction}

For single-step-ahead prediction, the LSTM is used in a \emph{many-to-one} configuration: an input sequence of $n_{\text{in}}$ time steps is processed by the LSTM layer(s), and only the hidden state at the final time step is passed through a fully connected (FC) output layer to produce the prediction \cite{goodfellow2016deep}. The complete architecture is:
\begin{center}
Input $(n_{\text{in}}, d)$ $\;\to\;$ LSTM layer(s) $\;\to\;$ last time step $\mathbf{h}_T$ $\;\to\;$ FC$(h \to d_{\text{out}})$ $\;\to\;$ Output
\end{center}

For multi-layer LSTMs ($L > 1$), each layer receives the hidden state sequence of the previous layer as input, enabling the network to learn hierarchical temporal representations at different time scales \cite{hermans2013training}. In this thesis, 1--2 layers are used depending on the parameter budget.

\subsubsection{Training}

LSTM networks are trained by minimizing a loss function---in this thesis, the mean squared error (MSE)---using gradient descent. Each training component is chosen based on established best practices in the literature:

\begin{itemize}
    \item \textbf{Adam optimizer} \cite{kingma2015adam}: Adapts the learning rate for each parameter based on estimates of the first and second moments of the gradient, making training more robust to the initial learning rate compared to standard stochastic gradient descent.

    \item \textbf{Mini-batch training} (batch size 32): Small batch sizes provide a regularizing effect through gradient noise and have been shown to generalize better than large batches \cite{masters2018revisiting}.

    \item \textbf{Dropout} \cite{zaremba2014recurrent}: Applied between LSTM layers (not within recurrent connections), which preserves long-term memory while regularizing effectively. Zaremba et al. showed that applying dropout only to non-recurrent connections substantially reduces overfitting. For single-layer LSTMs, inter-layer dropout is not applicable and is set to zero.

    \item \textbf{Early stopping} \cite{prechelt1998early}: Validation loss is monitored during training, and the model weights corresponding to the lowest validation loss are retained. Training halts when the validation loss has not improved for a specified number of epochs (patience). This acts as a form of implicit regularization, preventing the model from memorizing the training data.
\end{itemize}

\subsubsection{Strengths and Limitations}

LSTM networks have several advantages: they can learn complex nonlinear mappings from data, their gating mechanism effectively addresses the vanishing gradient problem, and they scale well with modern GPU hardware. However, they also have notable limitations: training requires large datasets and many gradient updates, the computational cost scales with the sequence length and number of parameters, and the model is a black box that offers limited interpretability.

% ============================================================================
\subsection{Reservoir Computing and Echo-State Networks}
\label{sec:esn}
% ============================================================================

Reservoir Computing (RC) is a computational framework that takes a fundamentally different approach to processing temporal data \cite{lukosevicius2009reservoir, tanaka2019recent}. Rather than training all network weights through gradient descent, RC exploits the dynamics of a fixed, randomly initialized \emph{reservoir} of recurrently connected units. Only the \emph{readout layer} (output weights) is trained, typically using simple linear regression. This separation between the dynamic reservoir and the trainable readout makes RC extremely fast to train.

The concept emerged independently in two forms: Echo-State Networks (ESNs) introduced by Jaeger \cite{jaeger2001echo} and Liquid State Machines introduced by Maass et al. \cite{maass2002real}. This thesis focuses on ESNs, which operate on discrete-time signals and are more widely used for time-series prediction.

\subsubsection{ESN Architecture}

An Echo-State Network consists of three components:

\begin{enumerate}
    \item \textbf{Input layer}: An input weight matrix $\mathbf{W}_{\text{in}} \in \mathbb{R}^{N \times d}$ maps the $d$-dimensional input signal to the $N$-dimensional reservoir space.
    \item \textbf{Reservoir}: A recurrent weight matrix $\mathbf{W} \in \mathbb{R}^{N \times N}$ defines the internal connections of $N$ reservoir nodes. This matrix is randomly generated (typically sparse) and kept fixed after initialization.
    \item \textbf{Readout}: An output weight matrix $\mathbf{W}_{\text{out}} \in \mathbb{R}^{d_{\text{out}} \times N}$ maps the reservoir states to the output. This is the only trained component.
\end{enumerate}

The reservoir state $\mathbf{r}_t \in \mathbb{R}^N$ is updated as:
\begin{equation}
    \mathbf{r}_t = (1 - \alpha) \, \mathbf{r}_{t-1} + \alpha \, \tanh(\mathbf{W} \, \mathbf{r}_{t-1} + \mathbf{W}_{\text{in}} \, \mathbf{x}_t),
    \label{eq:esn_state}
\end{equation}
where $\alpha \in (0, 1]$ is the \emph{leakage rate}, which controls the speed of the reservoir dynamics. When $\alpha = 1$, the update reduces to a standard RNN update; smaller values introduce exponential smoothing, slowing the dynamics.

The output is computed as:
\begin{equation}
    \mathbf{y}_t = \mathbf{W}_{\text{out}} \, \mathbf{r}_t.
    \label{eq:esn_output}
\end{equation}

\subsubsection{The Echo State Property}

For an ESN to function correctly, it must satisfy the \emph{Echo State Property} (ESP): the effect of initial reservoir states must fade over time, so that the current reservoir state depends only on the input history \cite{jaeger2001echo}. A sufficient condition for the ESP is that the \emph{spectral radius} $\rho(\mathbf{W})$---the largest absolute eigenvalue of the reservoir weight matrix---satisfies $\rho(\mathbf{W}) < 1$.

In practice, the spectral radius is a critical hyperparameter:
\begin{itemize}
    \item Small spectral radius ($\rho \ll 1$): Fast-decaying reservoir dynamics, short memory, suitable for tasks requiring rapid adaptation.
    \item Spectral radius near 1 ($\rho \lesssim 1$): Slow-decaying dynamics, longer memory, suitable for tasks with long-range dependencies.
    \item Spectral radius above 1 ($\rho > 1$): The ESP may be violated in principle, though input-driven operation can stabilize the dynamics. Values slightly above 1 sometimes yield good results for chaotic systems \cite{lukosevicius2012practical}.
\end{itemize}

\subsubsection{Training via Ridge Regression}

The readout weights are computed by solving a linear regression problem. Given the collected reservoir states $\mathbf{R} = [\mathbf{r}_1, \ldots, \mathbf{r}_T] \in \mathbb{R}^{N \times T}$ and target outputs $\mathbf{Y} = [\mathbf{y}_1, \ldots, \mathbf{y}_T] \in \mathbb{R}^{d_{\text{out}} \times T}$, the readout weights are obtained via Tikhonov (ridge) regression \cite{tikhonov1963solution}:
\begin{equation}
    \mathbf{W}_{\text{out}} = \mathbf{Y} \mathbf{R}^\top (\mathbf{R} \mathbf{R}^\top + \beta \mathbf{I})^{-1},
    \label{eq:ridge}
\end{equation}
where $\beta > 0$ is the ridge (regularization) parameter and $\mathbf{I}$ is the identity matrix. This closed-form solution avoids iterative optimization entirely: no gradients, no learning rate, no epochs. The entire training reduces to a single matrix inversion, which is extremely fast for moderate reservoir sizes.

\subsubsection{Key Hyperparameters}

ESN performance is sensitive to several hyperparameters \cite{lukosevicius2012practical, matzner2022hyperparameter}:
\begin{itemize}
    \item \textbf{Spectral radius} ($\rho$): Controls the reservoir's memory and dynamic range.
    \item \textbf{Leakage rate} ($\alpha$): Controls the speed of reservoir dynamics.
    \item \textbf{Reservoir size} ($N$): The number of reservoir nodes. Larger reservoirs provide more computational capacity but increase the number of total parameters.
    \item \textbf{Density} ($\delta$): The fraction of non-zero entries in $\mathbf{W}$. Sparse reservoirs (1--5\% connectivity) provide a form of implicit regularization by limiting inter-node coupling, which promotes diverse dynamic substructures within the reservoir and reduces overfitting. In practice, very sparse reservoirs often perform comparably to or better than dense ones while being computationally cheaper \cite{lukosevicius2009reservoir}.
    \item \textbf{Input fraction}: The fraction of reservoir nodes that receive input connections.
    \item \textbf{Ridge parameter} ($\beta$): Controls regularization strength; prevents overfitting to noisy training data.
\end{itemize}

Research has shown that ESN performance is highly sensitive to hyperparameter choices, particularly the spectral radius and leakage rate, with optimal performance occurring in a narrow region of the hyperparameter space called the ``edge of criticality'' \cite{matzner2022hyperparameter, mwamsojo2024stochastic}. This sensitivity necessitates systematic hyperparameter search, as discussed in Chapter~\ref{chapter3}.

\subsubsection{Parameter Count and Comparison with LSTM}
\label{sec:param_count}

A crucial aspect of comparing ESNs and LSTMs is the distinction between \emph{total} and \emph{trainable} parameters. For an ESN with $N$ reservoir nodes, input dimension $d$, density $\delta$, and input fraction $f_{\text{in}}$:
\begin{align}
    N_{\text{reservoir}} &= N^2 \times \delta, \label{eq:esn_reservoir} \\
    N_{\text{input}} &= N \times d \times f_{\text{in}}, \label{eq:esn_input} \\
    N_{\text{readout}} &= N \times d_{\text{out}}, \label{eq:esn_readout} \\
    N_{\text{total}} &= N_{\text{reservoir}} + N_{\text{input}} + N_{\text{readout}}, \label{eq:esn_total} \\
    N_{\text{trainable}} &= N_{\text{readout}}. \label{eq:esn_trainable}
\end{align}
Typically, only 3--5\% of the total parameters are trainable (the readout weights), while the vast majority (reservoir and input weights) are fixed after random initialization.

In contrast, an LSTM has $N_{\text{total}} = N_{\text{trainable}}$: every parameter is optimized through backpropagation. This fundamental difference raises the question of what constitutes a ``fair'' comparison: should models be matched by total parameters (which determines memory footprint and the dimension of the dynamical system) or by trainable parameters (which determines optimization difficulty)? In this thesis, we adopt the total parameter budget as the matching criterion, as it reflects the overall model capacity and computational resources. This choice is further discussed in Chapter~\ref{chapter5}.

\subsubsection{The pyReCo Library}

The ESN implementation in this thesis uses the \texttt{pyReCo} (Python Reservoir Computing) library, which provides a high-level interface for constructing, training, and evaluating ESNs. The library handles reservoir construction (random generation with specified spectral radius, density, and input fraction), state collection during the input-driven phase, and readout training via ridge regression. It supports time-series cross-validation for hyperparameter selection.

% ============================================================================
\subsection{Related Work}
\label{sec:related}
% ============================================================================

The comparison of Reservoir Computing with gradient-trained recurrent neural networks has attracted growing interest. This section reviews key prior work and identifies the gaps addressed by this thesis.

\subsubsection{RC for Chaotic Prediction}

Jaeger and Haas \cite{jaeger2004harnessing} demonstrated that ESNs can predict chaotic time series with remarkable accuracy, outperforming contemporary methods on the Mackey-Glass and Lorenz benchmarks. Pathak et al. \cite{pathak2018model} showed that RC can predict the short-term evolution of high-dimensional chaotic systems (Kuramoto-Sivashinsky equation) and reproduce their long-term statistical properties. These results established RC as a leading paradigm for chaotic prediction.

\subsubsection{Comparative Studies}

Vlachas et al. \cite{vlachas2020backpropagation} compared backpropagation-trained RNNs with RC for spatiotemporal chaotic systems, finding that both approaches achieve comparable short-term prediction accuracy but that RC better reproduces the long-term climate of the attractor. Chattopadhyay et al. \cite{chattopadhyay2020data} compared RC, artificial neural networks, and LSTM on the Lorenz 96 system, concluding that RC achieves the best trade-off between accuracy and computational cost.

Jirak et al. \cite{jirak2023echo} compared ESNs and LSTMs for gesture recognition, noting the lack of established benchmarking standards and the difficulty of making fair comparisons. Kumar and Kaushal \cite{kumar2025forecasting} applied RC to crude oil price forecasting, highlighting the computational efficiency advantage of RC over LSTM.

\subsubsection{Gaps in the Literature}

Despite these studies, several gaps remain:
\begin{enumerate}
    \item \textbf{Matched parameter budgets}: Most prior comparisons do not control for model size. A model with 50\,000 parameters is expected to outperform one with 1\,000 parameters, regardless of architecture. This thesis explicitly matches total parameter budgets across architectures.
    \item \textbf{Data efficiency}: Few studies systematically vary the amount of training data to assess how quickly each architecture learns. This thesis evaluates performance across six data lengths (1\,000--10\,000 samples).
    \item \textbf{Green computing}: Environmental cost metrics (training time, energy consumption, memory usage) are rarely reported alongside prediction accuracy. This thesis includes comprehensive efficiency metrics.
    \item \textbf{Statistical rigor}: Many comparisons report single-run results without confidence intervals or significance tests. This thesis uses five random seeds, paired t-tests, and effect size measures for all comparisons.
\end{enumerate}

A recent review of RC benchmarking practices \cite{wringe2025benchmarks} emphasized the need for standardized evaluation methodologies, consistent metrics, and reproducible experimental setups---principles that guide the methodology of this thesis.
